\chapter{Fundamentação Teórica}
\label{cha:fundamentacao}

\section{HIPÓTESE DE EFICIÊNCIA DE MERCADO}

A Hipótese de Eficiência de Mercado (\textit{Efficient Market Hypothesis} -- EMH), formalizada por \citep{Fama1970}, postula que os preços dos ativos financeiros refletem integralmente todas as informações disponíveis. Consequentemente, não seria possível obter retornos anormais de forma consistente utilizando exclusivamente informações públicas ou históricas, uma vez que tais dados já estariam incorporados aos preços.

No mesmo trabalho, Fama classifica a eficiência de mercado em três níveis:

\begin{itemize}
\item Forma fraca: os preços refletem apenas informações históricas de preço e volume;
\item Forma semi-forte: os preços incorporam todas as informações públicas disponíveis;
\item Forma forte: os preços refletem todas as informações, inclusive as privadas.
\end{itemize}

A relevância dessa hipótese para o presente trabalho reside no debate sobre a previsibilidade dos retornos financeiros. Em um mercado perfeitamente eficiente, estratégias quantitativas baseadas em padrões históricos não seriam capazes de gerar retornos superiores de forma consistente. Contudo, evidências empíricas apontam para a existência de anomalias de mercado, como os efeitos momentum e \textit{value premium} \citep{FamaFrench1992, JegadeeshTitman1993}.

\subsection{Críticas Comportamentais e a Hipótese de Mercados Adaptativos}

Apesar de sua relevância teórica, a EMH recebeu críticas substanciais da economia comportamental. Estudos de \citep{Shiller2000} e \citep{Thaler1993} demonstram que vieses cognitivos, como excesso de confiança, aversão à perda e comportamento de manada, podem levar a desvios persistentes entre preços de mercado e valores fundamentais.

Buscando reconciliar a racionalidade da EMH com as evidências comportamentais, \citet{Lo2004} propôs a Hipótese de Mercados Adaptativos (\textit{Adaptive Markets Hypothesis} -- AMH). Neste modelo, a eficiência de mercado não é um estado fixo, mas um processo dinâmico que evolui ao longo do tempo conforme os agentes aprendem e se adaptam às mudanças no ambiente econômico.

Essa perspectiva fornece suporte teórico para o uso de sistemas baseados em inteligência artificial e modelos de linguagem, uma vez que modelos adaptativos e multiagentes podem capturar padrões dinâmicos e interpretar contextos que estratégias estáticas ou puramente matemáticas não conseguem identificar.

\section{BENCHMARKS SIMPLES: BUY AND HOLD E EQUAL WEIGHT}

Estratégias quantitativas tradicionais e \textit{benchmarks} cumprem o papel de base comparativa essencial na avaliação de sistemas de decisão, pois oferecem referências de implementação simples e menos dependentes de estimações sofisticadas, o que facilita o contraste com abordagens mais complexas de alocação e seleção de sinais \citep{Markowitz1952, Clarke2006, LopezDePrado2018}.

\subsection{Buy and Hold}
\textit{Buy and Hold} é uma estratégia passiva em que o investidor compra um ativo ou uma carteira de ativos e permanece posicionado durante todo o período de avaliação. Em estudos de \textit{backtesting}, sua utilidade metodológica reside em fornecer uma linha de base de baixa complexidade, útil para verificar se uma estratégia ativa gera valor além da simples exposição ao mercado após a consideração de risco e custos operacionais \citep{LopezDePrado2018, Bailey2017}. No contexto deste trabalho, esse \textit{benchmark} permite distinguir eventuais ganhos associados ao processo de deliberação dos agentes daqueles que decorreriam apenas do comportamento geral dos ativos analisados.

\subsection{Equal Weight}
\textit{Equal Weight}, ou carteira igualmente ponderada, consiste em distribuir o capital igualmente entre os ativos selecionados. Sua relevância decorre do fato de incorporar diversificação sem recorrer diretamente à estimação de parâmetros como retorno esperado e matriz de covariância, exigidos por modelos de otimização clássicos \citep{Markowitz1952, Clarke2006, DeMiguel2009}. Por essa razão, trata-se de um \textit{benchmark} particularmente útil para comparar abordagens sofisticadas com uma regra simples de alocação, permitindo avaliar se o aumento de complexidade metodológica se traduz, de fato, em ganhos consistentes de desempenho ajustado ao risco.

\section{ESTRATÉGIAS BASEADAS EM ANÁLISE TÉCNICA}

A análise técnica assume que padrões históricos de preços e indicadores derivados dessas séries temporais podem gerar sinais úteis para previsão de movimentos futuros do mercado. \citet{Brock1992} investigam regras técnicas simples e mostram que tais procedimentos podem ser avaliados empiricamente por meio de simulações históricas. Da mesma forma, \citet{Lo2000} tratam padrões gráficos sob uma perspectiva estatística, reforçando o uso de estratégias técnicas como \textit{benchmarks} quantitativos.

\subsection{Estratégias de Tendência (Trend Following)}

Estratégias de \textit{trend following} buscam explorar a persistência de movimentos de preço. O \textit{Moving Average Crossover} é uma estratégia técnica clássica baseada no cruzamento entre médias móveis de diferentes horizontes. \citet{Brock1992} demonstram que regras baseadas em médias móveis apresentam poder preditivo estatisticamente significativo. A média móvel simples é definida como:

\begin{equation}
MA_t = \frac{1}{n} \sum_{i=0}^{n-1} P_{t-i}
\end{equation}

Quando uma média curta cruza acima de uma média longa, interpreta-se como um sinal de tendência positiva. Essa regra traduz a ideia de acompanhamento de tendência em um critério operacional claro.

Outra medida amplamente utilizada é o momentum, definido como o retorno acumulado em uma janela temporal:

\begin{equation}
Momentum_t = \frac{P_t}{P_{t-n}} - 1
\end{equation}

\citet{JegadeeshTitman1993} mostram que estratégias baseadas em momentum produzem retornos anormais positivos. A ideia central é que ativos com desempenho positivo em uma janela recente tendem a continuar apresentando desempenho superior \citep{Moskowitz2012}.

\subsection{Estratégias de Reversão à Média (Mean Reversion)}

A hipótese de reversão à média (\textit{Mean Reversion}) pressupõe que desvios expressivos em relação a um nível considerado normal tendem a ser parcialmente revertidos ao longo do tempo. \citet{PoterbaSummers1988} investigam componentes transitórios e encontram evidências de reversão à média em horizontes de longo prazo. O sinal de negociação pode ser obtido por meio do z-score:

\begin{equation}
z_t = \frac{P_t - \mu}{\sigma}
\end{equation}

onde $\mu$ representa a média móvel e $\sigma$ o desvio padrão. A lógica consiste em comprar ativos que registraram queda recente ou reduzir a exposição a ativos que subiram acentuadamente.

Entre os indicadores mais utilizados destacam-se o Índice de Força Relativa (RSI), proposto por \citet{Wilder1978}, e as Bandas de Bollinger, introduzidas por \citet{Bollinger2001}. As Bandas de Bollinger são definidas por:

\begin{equation}
BB_{upper} = MA + k\sigma
\end{equation}

\begin{equation}
BB_{lower} = MA - k\sigma
\end{equation}

O RSI busca medir a intensidade dos movimentos recentes de alta e baixa de um ativo, identificando sobrecompra ou sobrevenda. A esses indicadores soma-se o MACD (\textit{Moving Average Convergence Divergence}), que combina a diferença entre duas médias exponenciais e sua linha de sinal para identificar mudanças de tendência. A comparação entre \textit{momentum} e \textit{mean reversion} é relevante porque essas estratégias representam hipóteses opostas sobre a dinâmica dos retornos e, em mercados sujeitos a mudanças de regime, podem atuar de forma complementar \citep{JegadeeshTitman1993, PoterbaSummers1988, Moskowitz2012}.

\section{FACTOR INVESTING E MODELOS DE PRECIFICAÇÃO DE ATIVOS}

Enquanto a análise técnica foca em padrões históricos de preços, o \textit{Factor Investing} busca explicar retornos a partir de características fundamentais das empresas. O modelo de três fatores de \citep{FamaFrench1993} é um dos mais influentes:

\begin{equation}
R_i - R_f = \beta_m (R_m - R_f) + \beta_s SMB + \beta_v HML
\end{equation}

Posteriormente, \citep{FamaFrench2015} expandiram o modelo. Estudos recentes demonstram que modelos de \textit{machine learning} podem capturar relações não lineares entre fatores de risco e retornos esperados, superando abordagens lineares tradicionais \citep{Gu2020}.

\section{TEORIA MODERNA DE PORTFÓLIO}

A Teoria Moderna de Portfólio, proposta por \citet{Markowitz1952}, estabelece uma das bases mais importantes da teoria financeira moderna. A decisão de investimento não deve considerar apenas o retorno esperado de cada ativo isoladamente, mas também o risco da carteira como um todo e as relações de covariância entre os ativos. A diversificação permite reduzir o risco total do portfólio sem necessariamente reduzir o retorno esperado. O risco de um portfólio é definido pela variância:

\begin{equation}
\sigma_p^2 = w^T \Sigma w
\end{equation}

Essa abordagem oferece um \textit{benchmark} clássico de alocação, útil como referência para comparação com estratégias mais sofisticadas.

\subsection{Carteira de Mínima Variância}
A carteira de mínima variância é uma aplicação direta da estrutura de Markowitz. Seu objetivo é encontrar os pesos dos ativos que minimizam a variância total da carteira, dadas as estimativas de covariância. No contexto do mercado acionário brasileiro, pode funcionar como um \textit{benchmark} conservador \citep{Clarke2006}, concentrando-se na redução de risco quando as estimativas de retorno esperado são instáveis.

\subsection{Relação Risco-Retorno}
A comparação entre estratégias exige considerar conjuntamente retorno, risco e custos operacionais, de modo que o desempenho observado não seja interpretado apenas pelo lucro absoluto \citep{LopezDePrado2018}.

\section{SISTEMAS MULTIAGENTES E MODELOS DE LINGUAGEM EM FINANÇAS}

O uso de \textit{Large Language Models} (LLMs) em finanças quantitativas vem sendo investigado como uma alternativa para incorporar contexto, raciocínio textual e coordenação de tarefas em processos de decisão que tradicionalmente dependem de regras fixas ou modelos puramente numéricos. Em particular, o paradigma do \textit{Agentic Trading} descrito por \citet[p. 1]{Xia2026} trata esses modelos como agentes capazes de recuperar informação, estruturar argumentos e emitir decisões em ciclos iterativos. Ainda assim, a literatura recente sugere que o ganho potencial desses sistemas depende menos do uso isolado de um LLM e mais da forma como memória, especialização funcional e mecanismos de verificação são organizados no fluxo decisório.

\subsection{Arquitetura Multiagente e a Atribuição de Personas}

A literatura sobre sistemas multiagentes parte da hipótese de que a decomposição do processo decisório em papéis especializados pode reduzir ambiguidades e melhorar a consistência do raciocínio. Em vez de concentrar todas as tarefas em um único modelo, \citet[p. 1]{Li2023} mostram que a combinação de personas distintas, apoiadas por memória em camadas, pode melhorar a qualidade da decisão em cenários de \textit{trading}. Sob essa perspectiva, a atribuição de funções como análise técnica, gestão de risco e gestão de portfólio deixa de ser apenas uma escolha de implementação e passa a constituir um mecanismo de controle do escopo analítico de cada agente.

\subsection{O Fluxo de Debate Heterogêneo e Mitigação de Erros}

Outra linha relevante enfatiza mecanismos explícitos de debate e verificação entre agentes. Em \citet{Xing2024}, o modelo de \textit{Heterogeneous Agent Discussion} (HAD) organiza agentes com papéis analíticos distintos e orientações voltadas à identificação de erros específicos, buscando reduzir incertezas gerativas e vieses do processo inferencial. O interesse dessa abordagem, para o presente trabalho, não está em assumir que o debate multiagente substitui automaticamente outras técnicas, mas em reconhecer que a heterogeneidade de papéis pode funcionar como camada adicional de validação quando comparada a inferências isoladas do tipo \textit{Zero-shot} \citep[p. 5, p. 9]{Xing2024}.

\subsection{Estruturação do Raciocínio Encadeado (Chain-of-Thought)}

Além da especialização por papéis, a literatura recente destaca a importância de estruturar o próprio raciocínio produzido pelo modelo. Enquanto abordagens de \textit{Zero-shot Chain-of-Thought} permitem respostas mais livres, trabalhos como \citet{FinCoT2025} defendem que, em domínios financeiros, o alinhamento melhora quando o raciocínio é guiado por passos e critérios pré-definidos por especialistas. \citet[p. 1]{FinCoT2025} evidenciam que, ao fundamentar o processo gerativo em passos lógicos guiados por domínio (prática cunhada como FinCoT), a acurácia de alinhamento financeiro de modelos gerais de ponta se eleva drasticamente, com melhorias reportadas de 63,2\% para 80,5\%. Para este projeto, essa discussão é relevante porque sugere que a utilidade de uma arquitetura multiagente não depende apenas da presença de vários agentes, mas também do grau de estrutura imposto ao fluxo de análise, verificação e decisão.

\section{BACKTESTING EM SÉRIES FINANCEIRAS E RISCOS METODOLÓGICOS}

\textit{Backtesting} é o processo de simular uma estratégia de investimento sobre dados históricos, aplicando regras de decisão no passado. Estratégias quantitativas estão sujeitas ao risco de \textit{overfitting}, onde resultados positivos decorrem de seleções excessivamente ajustadas ou de vazamento de informação futura.

\citet{White2000} discute o problema de \textit{data snooping}, que ocorre quando muitas estratégias ou parâmetros são testados sobre o mesmo conjunto de dados e apenas os melhores resultados são reportados. \citet{Hansen2005} complementa abordando o teste de habilidade preditiva superior para evitar que estratégias lucrativas sejam resultado puramente estatístico. Adicionalmente, \citet{Bailey2017} demonstram que a probabilidade de encontrar estratégias aparentemente lucrativas (\textit{Probability of Backtest Overfitting}) aumenta significativamente quando múltiplos modelos são testados simultaneamente.

\section{MÉTRICAS DE DESEMPENHO E RISCO}

A comparação entre estratégias de investimento exige métricas capazes de capturar diferentes dimensões, mitigando análises enviesadas focadas estritamente no lucro absoluto.

\subsection{Retorno Acumulado e Sharpe Ratio}
O retorno acumulado mede a variação total do capital (\(R_{acum} = \prod_{t=1}^{T} (1 + R_t) - 1\)). Contudo, deve ser interpretado com o risco assumido. O \textit{Sharpe ratio} \citep{Sharpe1994, Lo2002} resume a relação entre retorno excedente e a volatilidade total (desvio-padrão), sendo definido como:

\begin{equation}
S = \frac{E[R_p - R_f]}{\sigma_p}
\end{equation}

\subsection{Sortino Ratio e Maximum Drawdown}
O \textit{Sortino ratio} \citep{Sortino1994} considera apenas o risco de queda (\textit{downside deviation}), avaliando o retorno excedente em relação a perdas abaixo de um limite aceitável (MAR). O \textit{Maximum Drawdown} \citep{Magdon2004} captura a severidade da maior perda acumulada entre um pico e um vale. Ambas as métricas são cruciais para avaliar a estabilidade do sistema de agentes frente a crises do mercado brasileiro.

\subsection{Turnover e Custos de Transação}
O \textit{turnover} mede a intensidade de negociação de uma estratégia. Estratégias ativas que rebalanceiam posições com muita frequência incorrem em maiores despesas. \citet{Harris2003} discute aspectos de microestrutura e custos de transação (corretagem, spread e emolumentos). A avaliação exige que o retorno líquido após a inclusão desses custos justifique a viabilidade da aplicação prática.
